% ===== PRÉAMBULE CANONIQUE — à coller VERBATIM en tête de CHAQUE .tex =====
\documentclass[11pt,a4paper]{article}
\usepackage[utf8]{inputenc}\usepackage[T1]{fontenc}\usepackage{lmodern}
\usepackage[french]{babel}
\usepackage{amsmath,amssymb,mathtools}\usepackage{icomma}
\usepackage[version=4]{mhchem}          % \ce{...} : équations & formules
\usepackage{siunitx}\sisetup{locale=FR} % \qty{}{}, \si{}, \ang{}
\usepackage{chemfig}                    % \chemfig{...} : structures développées
\usepackage{xcolor}\usepackage{colortbl}
\usepackage{tikz}\usepackage{pgfplots}\pgfplotsset{compat=1.18}
\usepackage[most]{tcolorbox}\usepackage{enumitem}\usepackage{array,booktabs}
\usepackage[a4paper,margin=2.2cm]{geometry}
\usetikzlibrary{arrows.meta,calc,angles,quotes,positioning,patterns,backgrounds,shapes.geometric}
\definecolor{primary}{RGB}{222,49,99}\definecolor{primarydark}{RGB}{150,22,55}
\definecolor{defcol}{RGB}{0,114,178}\definecolor{methcol}{RGB}{52,92,148}
\definecolor{excol}{RGB}{0,135,90}\definecolor{attncol}{RGB}{200,40,55}
\tcbset{boxbase/.style={enhanced,breakable,boxrule=0.9pt,arc=2.5pt,
  left=3mm,right=3mm,top=2.3mm,bottom=2.3mm,fonttitle=\bfseries,coltitle=white}}
% --- boîtes TUTORIEL (noms EXACTS, ne pas renommer) ---
\newtcolorbox{cle}[1][]{boxbase,colback=defcol!6,colframe=defcol,
  title={\textsc{À retenir}\ifx\relax#1\relax\relax\else\textmd{ : \itshape #1}\fi}}
\newtcolorbox{methode}[1][]{boxbase,colback=methcol!7,colframe=methcol,sharp corners,
  title={\textsc{Méthode}\ifx\relax#1\relax\relax\else\textmd{ --- \itshape #1}\fi}}
\newtcolorbox{exemplebox}[1][]{boxbase,colback=excol!5,colframe=excol,
  title={\textsc{Exemple}\ifx\relax#1\relax\relax\else\textmd{ : \itshape #1}\fi}}
\newtcolorbox{piege}[1][]{boxbase,colback=attncol!6,colframe=attncol,
  title={\textsc{Piège}\ifx\relax#1\relax\relax\else\textmd{ : \itshape #1}\fi}}
\newtcolorbox{remarque}[1][]{boxbase,colback=black!4,colframe=black!45,
  title={\textsc{Remarque}\ifx\relax#1\relax\relax\else\textmd{ : \itshape #1}\fi}}
% --- boîtes EXERCICE / CORRIGÉ (noms EXACTS, auto-comptées) ---
\newtcolorbox[auto counter]{exo}[1][]{boxbase,colback=primary!4,colframe=primary,
  title={\textsc{Exercice}~\thetcbcounter\ifx\relax#1\relax\relax\else\textmd{ --- \itshape #1}\fi}}
\newtcolorbox[auto counter]{corrige}[1][]{boxbase,colback=excol!4,colframe=excol,
  title={\textsc{Corrigé}~\thetcbcounter\ifx\relax#1\relax\relax\else\textmd{ --- \itshape #1}\fi}}
\newcommand{\R}{\mathbb{R}}\newcommand{\N}{\mathbb{N}}\newcommand{\Z}{\mathbb{Z}}\newcommand{\ds}{\displaystyle}
% (un \usepackage{fancyhdr}+en-tête est possible ; il est ignoré côté web)
% =========================================================================
\begin{document}
\sisetup{per-mode=symbol}
\begin{corrige}[du sulfure de fer : partir de masses]
\begin{enumerate}[label=\alph*)]
  \item $n(\ce{Fe})=\dfrac{5,6}{56}=\qty{0.1}{\mole}$ ; $n(\ce{S})=\dfrac{6,4}{32}=\qty{0.2}{\mole}.$
  \item \ce{Fe} : $\frac{0,1}{1}=0,1$ ; \ce{S} : $\frac{0,2}{1}=0,2$. \textbf{\ce{Fe} est limitant}.
  \item $n(\ce{FeS})=n(\ce{Fe})=\qty{0.1}{\mole}$, d'où $m=0,1\times88=\qty{8.8}{\gram}.$
\end{enumerate}
\end{corrige}

\begin{corrige}[limitant en phase gazeuse]
\begin{enumerate}[label=\alph*)]
  \item \ce{CO} : $\frac{4}{2}=2$ ; \ce{O2} : $\frac{3}{1}=3$. \textbf{\ce{CO} est limitant}.
  \item $V(\ce{CO2})=V(\ce{CO})\times\dfrac{2}{2}=\qty{4}{\litre}.$
  \item \ce{O2} consommé $=4\times\frac{1}{2}=\qty{2}{\litre}$, donc restant $=3-2=\qty{1}{\litre}.$
\end{enumerate}
\end{corrige}

\begin{corrige}[attaque du zinc par l'acide]
\begin{enumerate}[label=\alph*)]
  \item $n(\ce{Zn})=\dfrac{6,5}{65}=\qty{0.1}{\mole}$ ; $n(\ce{HCl})=C\,V=1\times0,100=\qty{0.1}{\mole}.$
  \item \ce{Zn} : $\frac{0,1}{1}=0,1$ ; \ce{HCl} : $\frac{0,1}{2}=0,05$. \textbf{\ce{HCl} est limitant}.
  \item $n(\ce{H2})=n(\ce{HCl})\times\dfrac{1}{2}=0,1\times0,5=\qty{0.05}{\mole}$, d'où
        $V=0,05\times22,4=\qty{1.12}{\litre}.$
\end{enumerate}
\end{corrige}

\begin{corrige}[avancement maximal et bilan]
\begin{enumerate}[label=\alph*)]
  \item $\dfrac{n_0(\ce{A})}{2}=\dfrac{0,10}{2}=0,05$ ; $\dfrac{n_0(\ce{B})}{3}=\dfrac{0,20}{3}\approx0,067$.
        Donc $\xi_{\max}=\qty{0.05}{\mole}.$
  \item Le minimum est atteint par \ce{A} : \textbf{\ce{A} est limitant}.
  \item $n(\ce{A})=0,10-2(0,05)=\qty{0}{\mole}$ ; $n(\ce{B})=0,20-3(0,05)=\qty{0.05}{\mole}$ ;
        $n(\ce{C})=0,05=\qty{0.05}{\mole}$ ; $n(\ce{D})=6(0,05)=\qty{0.30}{\mole}.$
\end{enumerate}
\end{corrige}

\begin{corrige}[précipitation du sulfate de baryum]
\begin{enumerate}[label=\alph*)]
  \item $n(\ce{BaCl2})=0,5\times0,200=\qty{0.1}{\mole}$ ; $n(\ce{Na2SO4})=0,5\times0,100=\qty{0.05}{\mole}.$
  \item \ce{BaCl2} : $\frac{0,1}{1}=0,1$ ; \ce{Na2SO4} : $\frac{0,05}{1}=0,05$.
        \textbf{\ce{Na2SO4} est limitant}.
  \item $n(\ce{BaSO4})=n(\ce{Na2SO4})=\qty{0.05}{\mole}$, d'où $m=0,05\times233=\qty{11.65}{\gram}$ ;
        \ce{BaCl2} restant $=0,1-0,05=\qty{0.05}{\mole}.$
\end{enumerate}
\end{corrige}

\end{document}
